\section{\label{sec:summary}总结}

准分布与赝分布为核子结构打开了一扇新的窗口，并首次提供了直接从格点 QCD 提取部分子分布函数的机会。非微扰计算的初步结果令人鼓舞，但在格点学界能够宣称稳健地确定光前 PDF 之前，仍有挑战需要克服。格点 QCD 的一个特别困难，是与 Wilson 线长度相联系的幂次发散的存在，在取连续极限之前必须将其非微扰地移除。在文献~\cite{Monahan:2016bvm} 中，我们引入涂抹准分布作为规避 Wilson 线幂次发散问题的工具。利用梯度流——原始自由度在新参数流时下的经典演化——的重正化性质，涂抹准分布在连续极限下保持有限。

这里，我在单圈微扰论中研究了涂抹准分布。虽然梯度流通过引入零流时下不存在的新图并使相应积分复杂化而增加了微扰计算的难度，但其优势是显著的：只要把流时以物理单位固定，涂抹准分布在连续极限下就是有限的。所得的连续涂抹准分布随后可以直接、或经由零流时准分布，与 $\ms$ 方案中的光前 PDF 联系起来。准分布的红外行为不受流时影响，流时只用于调节紫外行为，因此匹配程序可以微扰地进行，即在外部规范固定的夸克态之间计算 Wilson 线算符。

我在固定 Wilson 线长度下进行匹配，此时外部夸克动量可以设为零。这减少了给定微扰阶下有贡献的图的数目，并简化了相应的单圈积分，后者可以解析完成。我用维数正规化调节对数红外发散，并证明由此产生的时空维数极点在涂抹与未涂抹准分布的匹配关系中抵消。我证明 Wilson 线长度趋于零的极限是良好定义的，并约化为有限流时下已知的矢量流结果。在小流时极限，矩阵元只对 Wilson 线长度与梯度流涂抹半径之比对数依赖。因此涂抹矩阵元满足一个与典型重正化群方程类似的方程，这为定义非微扰步标度程序提供了可能。最后，我通过数值计算相应的单圈图，研究了有限外部动量下的涂抹矩阵元。把这一研究扩展到两圈很可能需要自动化的方法，因为梯度流诱导的额外顶点使图的数目随微扰阶迅速增长。与涂抹准分布和赝分布相联系的系统不确定度的非微扰研究正在进行中。
